\Abstract
\textit{Oblivious Random Access Memory (ORAM)} enables secure access to storage, primarily in the context of databases, by obfuscating access patterns.
An attacker cannot infer any information from these access patterns.
\textit{Distributed ORAM (DORAM)} considers the partitioning of storage among multiple parties such that no single party can reconstruct the data by itself.
A fixed number of cooperating parties is required to reconstruct the secret, which means the stored data.
This research focuses on \textit{DUORAM}, a \textit{DORAM} implementation by Vadapalli et al.
While \textit{DUORAM} supports efficient read and write operations for 64-bit values, scaling it to arbitrary element sizes remains a challenge.
This research discusses and implements possible extensions of \textit{DUORAM} to support elements of arbitrary size.
Various approaches to scaling are discussed, along with all necessary technical foundations.
The proposed approaches have been implemented and benchmarked.
Tests regarding efficiency in terms of bandwidth, communication overhead, and storage requirements are executed.
The results provide insights into the trade-offs between computational complexity and practical feasibility, highlighting key limitations and potential optimizations for future research.

\section*{Abstract}
\textit{Oblivious Random Access Memory (ORAM)} ermöglicht sicheren Zugriff auf Speicher, meist im Kontext von Datenbanken, indem Zugriffssequenzen verborgen werden.
Einem Angreifer ist es nicht möglich Rückschlüsse auf die gespeicherten Daten oder zugrundeliegenden Datenstrukturen aus den Zugriffssequenzen zu ziehen.
\textit{Distributed ORAM (DORAM)} betrachtet die Aufteilung eines Speichers zwischen mehreren Parteien, sodass keine Partei alleine die Daten rekonstruieren kann.
Es benötigt eine feste Anzahl kooperierender Parteien um das Geheimnis, beziehungsweiße die Daten, wiederherzustellen.
Diese Research betrachtet \textit{DUORAM}, eine \textit{DORAM} Implementierung von Vadapalli et al.
Während \textit{DUORAM} effiziente Lese- und Schreiboperationen für 64 bit Werte unterstützt, bleibt die Skalierung auf beliebige Elementgrößen eine Herausforderung.
Diese Research erweitert und implementiert \textit{DUORAM}, um Elemente beliebiger Größe zu unterstützen.
Es werden verschiedene Ansätze zur Skalierung diskutiert, einschließich aller benötigten technischen Grundlagen.
Die Ansätze wurden implementiert und gebenchmarked.
Tests bezüglich Effizienz in Bezug auf Bandbreite, Kommunikationseffizienz und Speicherbedarf werden analysiert.
Die Ergebnisse liefern Einblicke in die Abwägungen zwischen Rechenkomplexität und praktischer Umsetzbarkeit und zeigen wesentliche Einschränkungen sowie mögliche weitere Optimierungen für künftige Forschung auf.